\documentclass[11pt,a4paper]{article}
\usepackage[margin=25mm]{geometry}
\usepackage{amsmath,amssymb,amsthm,booktabs}
\newtheorem{lemma}{Lemma}\newtheorem{theorem}{Theorem}
\title{Three order five counts from centralizers and generated subgroups}
\author{Human Compression project}\date{}
\begin{document}\maketitle
Put $G=O'N$, $g=|G|=460815505920$, and write
\[
 x_T(p,q)=|\{(r,s)\in T^2:|r|=p,\ |s|=q,\ |rs|=2\}|.
\]
The number of these pairs generating $T$ is denoted by $e_T(p,q)$.
We recover the three remaining entries involving order five in the
fixed Connor--Leemans table. The complete subgroup roster, the
primary ATLAS structures and character values are borrowed inputs.
Their exact editions remain bound to the accompanying readers.

For an actual subgroup $L\leq G$, its generating pairs are disjoint
from those of every other actual subgroup. Thus a conjugacy class
of subgroups of type $L$ contributes exactly
\[
                       [G:N_G(L)]e_L(p,q).
\tag{1}
\]
No containment multiplicities are needed in this identity. We first
remove two large families by a common centralizer argument. The
remaining generated subgroup types are small enough to treat by
literal permutations and elementary subgroup arguments.

\section*{The two distinguished maximal families}
We use row numbers from the complete roster. Let $\mathcal A$ be the
row 7 involution centralizers, and let $\mathcal B$ be the row 11
maximal subgroups. An accepted incidence argument proves that every
cyclic subgroup of order five lies in exactly one member of
$\mathcal B$. A member $B$ has a normal subgroup
\[
 B_0=F\times A_6,\qquad F=C_3^2:C_4,
\]
of index two. We also reuse the accepted row 32 intersection
subgroup $H$: it is maximal in its unique member of $\mathcal A$,
lies in a unique member of $\mathcal B$, and is their intersection.
It has a normal subgroup $N\cong A_6$ with $H/N\cong Q_8$.
The central subgroup of this quaternion quotient lifts to
$N\times C_2$. These structural statements were proved in the
earlier order 20 and order ten components.

\begin{lemma}\label{central}
For a pair with an entry of order five, outside all members of
$\mathcal B$, at most one member of $\mathcal A$ contains the pair.
For a fixed cyclic subgroup $C$ of order five, there are nine
members of $\mathcal A$ and nine row 32 subgroups containing $C$.
They correspond bijectively by intersection with the unique
$B\in\mathcal B$ containing $C$.
\end{lemma}
\begin{proof}
The primary character table has one element class of each of orders
three and five, with centralizer orders 3240 and 180. Choose a
nonidentity $c\in C_3^2\leq F$. The group $C_3^2\times A_6$, of
order 3240, centralizes $c$, so equals $C_G(c)$ and lies in $B$.
Likewise, for an element $d$ of order five in the $A_6$ factor,
$F\times\langle d\rangle$ has order 180 and equals $C_G(d)$.
Conjugacy gives the same containment in a member of $\mathcal B$
for every centralizer of an element of order three or five.

Let $r$ have order five, and let $C_G(r)\cap C_G(s)$ be the common
centralizer of a pair outside $\mathcal B$. It is a subgroup of a
group of order 180. If it had an element of order three or five,
both entries would lie in that element's centralizer, hence in a
member of $\mathcal B$, a contradiction. By Cauchy's theorem it
is therefore a 2-group. A Sylow 2-subgroup of
$C_G(r)=F\times C_5$ is cyclic of order four. The common
centralizer has at most one involution. Since every member of
$\mathcal A$ is the centralizer of an involution, the first claim
follows.

There is one cyclic order five class, of length 640021536. The
acquired trivial character counts 8064 elements of order five in
a row 7 subgroup, hence 2016 cyclic subgroups. Every row 32
subgroup has exactly 36 cyclic subgroups of order five: all lie
in its normal $A_6$, since its quotient has order eight. The direct
incidences are
\[
 \frac{2857239\cdot2016}{640021536}=9,
 \qquad
 \frac{160005384\cdot36}{640021536}=9.
\]
Each of the nine row 32 subgroups has unique maximal overgroups
in the two families. Its row 11 overgroup must be the unique $B$
containing $C$. Two of these row 32 subgroups cannot have the same
row 7 overgroup: each would then equal the intersection of that
overgroup with $B$. The nine-to-nine correspondence is therefore
bijective, as required.
\end{proof}

\begin{lemma}\label{covered}
The number $U(p,q)$ of pairs contained in the union of the two
distinguished families, for our three types, is
\[
 U(p,q)=2857239x_A(p,q)+17778376x_B(p,q)
                          -160005384x_H(p,q),
\tag{2}
\]
where
\[
 \begin{array}{c|rrr}
 (p,q)&(4,5)&(5,5)&(5,6)\\\hline
 x_H(p,q)&2880&2880&1440.
 \end{array}
\]
\end{lemma}
\begin{proof}
A pair belongs to at most one member of $\mathcal B$. Each row 7
overgroup of a pair already in $B$ gives exactly one row 32
intersection, by Lemma~\ref{central}. Thus the intersection term
in (2) cancels every extra row 7 contribution to such a pair.
Outside $\mathcal B$ a pair has at most one row 7 overgroup, again
by that lemma. Consequently every pair in the union is counted
exactly once in (2).

For the intrinsic values, an element of order five in $H$ lies in
$N=A_6$. If its product with another entry is an involution, the
other entry has image of order at most two in $H/N=Q_8$ and hence
lies in $A_6\times C_2$. An entry of order four has an order four
first coordinate and either sign. An entry of order six has an
order three first coordinate and negative sign, since $A_6$ has
no element of order six. Two entries of order five both lie in
$A_6$. The retained literal values
\[
 x_{A_6}(4,5)=1440,\quad x_{A_6}(5,5)=2880,
 \quad x_{A_6}(3,5)=1440
\]
give the three numbers. The accompanying reader checks them again
in the complete group $A_6\times C_2$.
\end{proof}

\section*{Small intrinsic counts}
The following table records generating pairs, rather than all
pairs. Zero means there is no generating pair of that type.
\[
\begin{array}{c|rrr}
 T&(4,5)&(5,5)&(5,6)\\\hline
 A_5&0&120&0\\
 S_5&120&0&120\\
 A_6&1440&1440&0\\
 L_2(11)&0&2640&2640\\
 L_2(31)&59520&59520&0\\
 J_1&0&351120&2106720
\end{array}
\tag{3}
\]

\begin{lemma}\label{small}
All the entries of (3) have the stated values.
\end{lemma}
\begin{proof}
The literal alternating permutations give
$x_{A_5}(5,5)=120$ and the three $A_6$ values above. Every counted
$A_5$ pair generates $A_5$: the proper maximal subgroups containing
five-elements are dihedral of order ten, and two rotations cannot
have product of order two. The complete maximal list of $A_6$
contains twelve $A_5$ subgroups; its other maximal orders are 36
and 24. Thus all $(4,5)$ pairs generate $A_6$, whereas
$e_{A_6}(5,5)=2880-12\cdot120=1440$.

The complete permutations of five letters give 120 pairs of each
of our three types in $S_5$. The $(5,5)$ pairs lie in its normal
$A_5$. A proper maximal subgroup can contain a $(4,5)$ pair only
if it is the affine group $C_5:C_4$; projecting that group onto
$C_4$ rules out a product of order two. No proper maximal subgroup
can contain a $(5,6)$ pair. Hence the other two counts are
generating counts. The reader independently closes every such
pair and verifies the order 120 result. For the roster's order
120 subgroups labelled $A_5:2$, their normal $A_5$ and the absence
of elements of order ten identify the group as $S_5$: a nontrivial
centralizer of that $A_5$ would give $A_5\times C_2$, while a
trivial centralizer embeds the whole group in
$\operatorname{Aut}(A_5)=S_5$.

For $L_2(11)$, a new literal model uses the determinant-one
fractional linear transformations on $\mathbb P^1(\mathbb F_{11})$.
The permutations $x\mapsto x+1$ and $x\mapsto-1/x$ generate all
660 transformations: upper and lower elementary matrices generate
$\operatorname{SL}_2(11)$ by elimination, and its kernel on the
projective line consists of the two scalar matrices. The reader
also compares this generated set with all determinant-one matrix
actions directly. The all-pair counts are $0,5280,2640$.
The complete maximal list has 22 subgroups $A_5$, twelve subgroups
of order 55, and 55 dihedral subgroups of order twelve. Only the
$A_5$ subgroups contribute proper $(5,5)$ pairs, and no proper
maximal contains a $(5,6)$ pair. Subtraction gives 2640 in both
cases.

Each $L_2(31)$ maximal list has two classes of 248 subgroups $A_5$,
with remaining maximal types of orders 465, 32, 30 and 24.
The odd order subgroup contributes no pair with product an
involution; in the dihedral subgroup of order thirty two
five-elements are rotations and again have no such product.
The other proper types exclude our orders. The acquired character
counts $59520,119040,0$ therefore give
$e_{L_2(31)}(4,5)=59520$ and
$e_{L_2(31)}(5,5)=119040-496\cdot120=59520$.

Finally use the complete maximal list of $J_1$. For $(5,6)$ its
proper generating subgroups are the 266 subgroups $L_2(11)$ and
the 1463 subgroups $A_5\times C_2$. In the latter group the count
is 120, with every pair generating the whole group. The type
$D_{10}\times D_6$ contributes zero: the order five entry lies in
the first rotation subgroup, while an order six entry has a
nontrivial order three coordinate in the second factor. Its
product cannot be an involution. The other maximal orders exclude
the pair. Hence
\[
 e_{J_1}(5,6)=2984520-266\cdot2640-1463\cdot120=2106720.
\]
For $(5,5)$, every proper generated subgroup is $L_2(11)$ or $A_5$.
In particular, both five-elements in $C_{11}:C_{10}$ lie in its
odd order subgroup of index two. For the exact number of its $A_5$ subgroups, we reuse the complete
table of marks of $J_1$ already acquired and accepted in the
$L036$ treatment. Its three classes of subgroups of order sixty
have normalizer orders $120,60,60$, respectively. The acquired
cyclic-class fusion and marks show that only the last contains
an element of order fifteen. The complete global roster permits
only $A_5$ and $D_{10}\times D_6$ as order sixty subgroups of
$J_1$; the latter contains fifteen-elements and the former does
not. Therefore the first two classes are exactly its $A_5$
classes, with lengths $175560/120=1463$ and $175560/60=2926$.
This is an explicit additional use of the previously acquired
complete mark table, rather than an assumption that a global
conjugacy class remains one class inside $J_1$. Thus
\[
 e_{J_1}(5,5)=1580040-266\cdot2640-4389\cdot120=351120.
\]
The zero $(4,5)$ entry follows from the acquired all-pair count.
All subtractions use exact generated subgroups, so their terms
are disjoint.
\end{proof}

\section*{Covering the $M_{11}$ and $A_7$ contributions}
\begin{lemma}\label{saturation}
There are no generating pairs of our three types in $M_{11}$ or
$A_7$. Every such pair is accounted for by the actual smaller
subgroups described below.
\end{lemma}
\begin{proof}
A row 23 subgroup $M_{11}$ has eleven maximal subgroups of row 102.
Each has a unique normal $A_6$ of row 168. These eleven $A_6$
subgroups are distinct, since their full normalizer order is 720
and equals their row 102 overgroup. There are also twelve maximal
$L_2(11)$ subgroups and 66 maximal $S_5$ subgroups of row 278.
The row 24 case has the identical statement with rows
103,169,279 instead.

Inside either $M_{11}$, a row 368 subgroup $A_5$ has normalizer
order 60, so its orbit has 132 members. Such a subgroup exists in
the indicated normal $A_6$. A row 370 subgroup $A_5$ exists as the
normal subgroup of one of the 66 maximal $S_5$ subgroups. Its
normalizer inside $M_{11}$ is exactly that $S_5$, because the
full normalizer has order 480 and the $S_5$ is maximal. Its orbit
has 66 members. The two orbits belong to different global rows,
so yield 198 distinct $A_5$ subgroups. By (3), disjoint generating
contributions now give
\[
\begin{array}{c|r|r}
(p,q)&\text{proper pairs exhibited}&x_{M_{11}}(p,q)\\\hline
(4,5)&11\cdot1440+66\cdot120&23760\\
(5,5)&11\cdot1440+12\cdot2640+198\cdot120&71280\\
(5,6)&12\cdot2640+66\cdot120&39600.
\end{array}
\]
The final column is the acquired character count, and equals the
middle column in every row. Thus these explicitly disjoint
contributions exhaust all pairs, with no need to infer any
unlisted subgroup incidence. This also rules out generating
pairs in the embedded row 102 and 103 subgroups of order 720.

A row 43 subgroup $A_7$ has seven maximal $A_6$ subgroups of row
163 and 21 maximal $S_5$ subgroups of row 280. The latter have
normal $A_5$ subgroups of row 369, giving an orbit of length 21:
their full normalizer is that same $S_5$, of order 120. There is
also a row 367 $A_5$ inside any of the indicated $A_6$ subgroups.
Its full normalizer is the row 281 $S_5$, which cannot lie in a
row 43 subgroup, as its complete overgroup list shows. Its
normalizer inside this $A_7$ is therefore just $A_5$, and its
orbit has length 42. These are 63 distinct $A_5$ subgroups. For
row 44, interchange 367 with 369 and 280 with 281. We obtain
\[
\begin{array}{c|r|r}
(p,q)&\text{proper pairs exhibited}&x_{A_7}(p,q)\\\hline
(4,5)&7\cdot1440+21\cdot120&12600\\
(5,5)&7\cdot1440+63\cdot120&17640\\
(5,6)&21\cdot120&2520.
\end{array}
\]
Again every exhibited pair has its specified exact generated
subgroup, and the sums equal the full character counts. This
proves the assertion.
\end{proof}

\section*{The affine contribution and the final count}
The accepted affine complement argument concerns the row 12
maximal subgroup $M=3^4:(2^{1+4}.D_{10})$. Its unique subgroup of
index two is row 20, of order 12960. All five-elements lie in
that subgroup. Its complete maximal list consists of an order
2592 subgroup, sixteen row 90 subgroups of order 810, and 81
row 245 subgroups $F=2^{1+4}:5$ of order 160. The order 810
subgroup has a normal subgroup of odd order 405 and index two,
so two five-elements cannot have an involution product there.
The order 2592 subgroup has no five-element.

The accepted extraspecial calculation has five nonzero singular
vectors in its four-dimensional quotient. The same calculation
gives $x_F(5,5)=64\cdot5=320$, equal to $x_F(10,10)$: multiplication
of both entries by the central involution preserves their product.
The 81 copies of $F$ give disjoint pair sets, because each
five-element has a unique fixed point in the affine action on
$\mathbb F_3^4$, hence belongs to a unique complement. Therefore
the exact generating count for row 20 is
\[
                       233280-81\cdot320=207360.
\tag{4}
\]
It has no pair of the other two types, since the containing
maximal subgroup's acquired character count is zero.

\begin{theorem}
The remaining ordered automorphism-orbit counts involving five are
\[
 \frac{e_G(4,5)}{2g}=18,\qquad
 \frac{e_G(5,5)}{2g}=26,\qquad
 \frac{e_G(5,6)}{2g}=98.
\]
\end{theorem}
\begin{proof}
The complete cyclic-order ancestor lists, after removing the
descendants of rows 7 and 11, leave the following possible
generated subgroups with nonzero counts. The last three columns
give their contributions in (1), divided by $g$, using (3) and (4).
\[
\begin{array}{c|c|rrr}
\text{type}&\text{global rows}&(4,5)&(5,5)&(5,6)\\\hline
J_1&6&0&2&12\\
L_2(31)&15,16&8&8&0\\
3^4:(2^{1+4}:5)&20&0&8&0\\
L_2(11)&119&0&4&4\\
A_6&163,168,169&8&8&0\\
S_5&280,281&2&0&2\\
A_5&367,368,369&0&4&0\\\hline
\text{total}&&18&34&18
\end{array}
\tag{5}
\]
For completeness about the discarded rows, $M_{11}$ and $A_7$
are disposed of by Lemma~\ref{saturation}; so are their indicated
subgroups of order 720. Both order five entries in an $S_5$ lie
in $A_5$, and the other two $S_5$ rows lie in the removed families.
Odd order groups cannot contain the product involution. The row
286 subgroup of order 110 has all five-elements in its odd order
subgroup of index two. The affine rows are treated by (4) and
the zero character counts. Every omitted row and its reason are
also checked in the finite support ledger.

Thus the number of proper pairs outside the distinguished union
is exactly $18g,34g,18g$. This uses the global normalizers of
the seven types in (5), counting each actual generated subgroup
once. Add $U(p,q)$ from (2) to obtain all nongenerating pairs,
and subtract from the acquired $x_G(p,q)$. Division by
$|\operatorname{Aut}(G)|=2g$ gives the three displayed values.
The automorphism action on generating pairs is free, and no
division by duality is made.
\end{proof}

\paragraph{Attribution and limits.}
The numerical assertions are from Connor and Leemans' fixed
Table 1; the present values agree with it. Complete subgroup
rosters, named ATLAS structures and character tables remain
borrowed premises. The common centralizer argument, exact
normalizer counts and saturation by disjoint generated subgroups
are the proof supplied here. The accepted earlier intersection
and affine-complement arguments are reused explicitly. The new
finite calculation constructs only $L_2(11)$ and $S_5$; no ambient
group census is used. This does not prove any outstanding
chirality assertion or claim Lean, whole-source or priority credit.
\end{document}
